\documentclass[10pt,a4paper]{amsart}
\usepackage{amsmath,amssymb,amsthm,graphicx}
\usepackage{verbatim}
\newenvironment{reference}{\comment}{\endcomment}
%\newenvironment{reference}{}{}
\newenvironment{slogan}{\comment}{\endcomment}
\newenvironment{history}{\comment}{\endcomment}
% For commutative diagrams we use Xy-pic
\usepackage[all]{xy}

% We use 2cell for 2-commutative diagrams.
\xyoption{2cell}
% Theorem environments.
%
\theoremstyle{plain}
\newtheorem{theorem}[subsection]{Theorem}
\newtheorem{proposition}[subsection]{Proposition}
\newtheorem{lemma}[subsection]{Lemma}

\theoremstyle{definition}
\newtheorem{definition}[subsection]{Definition}
\newtheorem{example}[subsection]{Example}
\newtheorem{exercise}[subsection]{Exercise}
\newtheorem{situation}[subsection]{Situation}

\theoremstyle{remark}
\newtheorem{remark}[subsection]{Remark}
\newtheorem{remarks}[subsection]{Remarks}

\numberwithin{equation}{subsection}

% Macros
%
\def\lim{\mathop{\mathrm{lim}}\nolimits}
\def\colim{\mathop{\mathrm{colim}}\nolimits}
\def\Spec{\mathop{\mathrm{Spec}}}
\def\Hom{\mathop{\mathrm{Hom}}\nolimits}
\def\Ext{\mathop{\mathrm{Ext}}\nolimits}
\def\SheafHom{\mathop{\mathcal{H}\!\mathit{om}}\nolimits}
\def\SheafExt{\mathop{\mathcal{E}\!\mathit{xt}}\nolimits}
\def\Sch{\mathit{Sch}}
\def\Mor{\mathop{\mathrm{Mor}}\nolimits}
\def\Ob{\mathop{\mathrm{Ob}}\nolimits}
\def\Sh{\mathop{\mathit{Sh}}\nolimits}
\def\NL{\mathop{N\!L}\nolimits}
\def\CH{\mathop{\mathrm{CH}}\nolimits}
\def\proetale{{pro\text{-}\acute{e}tale}}
\def\etale{{\acute{e}tale}}
\def\QCoh{\mathit{QCoh}}
\def\Ker{\mathop{\mathrm{Ker}}}
\def\Im{\mathop{\mathrm{Im}}}
\def\Coker{\mathop{\mathrm{Coker}}}
\def\Coim{\mathop{\mathrm{Coim}}}

% Boxtimes
%
\DeclareMathSymbol{\boxtimes}{\mathbin}{AMSa}{"02}

%
% Macros for moduli stacks/spaces
%
\def\QCohstack{\mathcal{QC}\!\mathit{oh}}
\def\Cohstack{\mathcal{C}\!\mathit{oh}}
\def\Spacesstack{\mathcal{S}\!\mathit{paces}}
\def\Quotfunctor{\mathrm{Quot}}
\def\Hilbfunctor{\mathrm{Hilb}}
\def\Curvesstack{\mathcal{C}\!\mathit{urves}}
\def\Polarizedstack{\mathcal{P}\!\mathit{olarized}}
\def\Complexesstack{\mathcal{C}\!\mathit{omplexes}}
% \Pic is the operator that assigns to X its picard group, usage \Pic(X)
% \Picardstack_{X/B} denotes the Picard stack of X over B
% \Picardfunctor_{X/B} denotes the Picard functor of X over B
\def\Pic{\mathop{\mathrm{Pic}}\nolimits}
\def\Picardstack{\mathcal{P}\!\mathit{ic}}
\def\Picardfunctor{\mathrm{Pic}}
\def\Deformationcategory{\mathcal{D}\!\mathit{ef}}
\usepackage{cleveref}
\crefname{originallemma}{original lemma}{original lemmas}
\Crefname{originallemma}{Original lemma}{Original lemmas}
\makeatletter
\newlabel{coherent-lemma-quasi-compact-h1-zero-covering}{{01XF}{source}{}{sourcepointer-1}{}}
\newlabel{coherent-lemma-quasi-compact-h1-zero-covering@cref}{{[originallemma][0][]01XF}{[1][0][]source}{}{}{}}
\makeatother
\crefname{originallemma}{original lemma}{original lemmas}
\Crefname{originallemma}{Original lemma}{Original lemmas}
\makeatletter
\newlabel{decent-spaces-lemma-there-is-a-scheme-integral-over}{{09YB}{source}{}{sourcepointer-2}{}}
\newlabel{decent-spaces-lemma-there-is-a-scheme-integral-over@cref}{{[originallemma][0][]09YB}{[1][0][]source}{}{}{}}
\makeatother
\crefname{originallemma}{original lemma}{original lemmas}
\Crefname{originallemma}{Original lemma}{Original lemmas}
\makeatletter
\newlabel{lemma-vanishing-base-change}{{07V0}{source}{}{sourcepointer-3}{}}
\newlabel{lemma-vanishing-base-change@cref}{{[originallemma][0][]07V0}{[1][0][]source}{}{}{}}
\makeatother
\crefname{originalreference}{original reference}{original references}
\Crefname{originalreference}{Original reference}{Original references}
\makeatletter
\newlabel{situation-vanishing}{{07UY}{source}{}{sourcepointer-4}{}}
\newlabel{situation-vanishing@cref}{{[originalreference][0][]07UY}{[1][0][]source}{}{}{}}
\makeatother
\crefname{originallemma}{original lemma}{original lemmas}
\Crefname{originallemma}{Original lemma}{Original lemmas}
\makeatletter
\newlabel{spaces-morphisms-lemma-image-universally-closed-separated}{{05Z2}{source}{}{sourcepointer-5}{}}
\newlabel{spaces-morphisms-lemma-image-universally-closed-separated@cref}{{[originallemma][0][]05Z2}{[1][0][]source}{}{}{}}
\makeatother
\crefname{originallemma}{original lemma}{original lemmas}
\Crefname{originallemma}{Original lemma}{Original lemmas}
\makeatletter
\newlabel{spaces-morphisms-lemma-integral-universally-closed}{{0415}{source}{}{sourcepointer-6}{}}
\newlabel{spaces-morphisms-lemma-integral-universally-closed@cref}{{[originallemma][0][]0415}{[1][0][]source}{}{}{}}
\makeatother
\title{Separatedness of the structure morphism in the vanishing situation, via an integral affine scheme cover}
\author{The Stacks Project Authors}
\date{stacks-project:spaces-cohomology.tex:local-snapshot-20261012; GNU FDL 1.2 or later}
\begin{document}
\maketitle
\noindent\textit{Editorial note.} All delivered mathematical text is the author's own, copied from the pinned Stacks Project chapter file spaces-cohomology.tex; no statement or proof step was rewritten, reordered or removed.\par\medskip
\noindent\textit{Editorial note.} Preamble provenance: the pinned chapter file opens with \textbackslash{}input\{preamble\} and carries no preamble text of its own, so the project's single Stacks-wide preamble staged at sources/stacks-spaces-cohomology/preamble.tex is excerpted verbatim as three spans, lines 8-12 (verbatim package and the project's own reference/slogan/history display environments), lines 14-18 (Xy-pic setup) and lines 154-217 (theorem declarations and mathematical macros). The wrapper class and the amsmath/amssymb/amsthm/graphicx base packages are supplied by the assembler. No macro was invented or redefined.\par\medskip
\noindent\textit{Editorial note.} Same-chapter pointers are located at their \textbackslash{}label lines inside the pinned chapter file; cross-chapter pointers are located in the project's own tag list staged in the snapshot (sources/stacks-chow/tags.txt), at the verbatim lines carrying their canonical tags.\par\medskip
\noindent\textit{Allowed standard prerequisites (editorial).}\begin{itemize}
\item Existence of a scheme integral and surjective over the space. Decent Spaces, label decent-spaces-lemma-there-is-a-scheme-integral-over (invoked at proof line 3139). It supplies the scheme \$Y\$ with a surjective integral morphism \$Y \textbackslash{}to X\$, which is affine and therefore covered by the base change vanishing.
\item Vanishing along an affine base change, and the affine criterion for a quasi-compact scheme. Same chapter spaces-cohomology.tex, label lemma-vanishing-base-change (invoked at proof line 3141) and Cohomology of Schemes, label coherent-lemma-quasi-compact-h1-zero-covering (invoked at proof line 3148). The first gives the vanishing of \$H\textasciicircum{}1\$ of quasi-coherent modules on \$Y\$; the second, since \$Y\$ is a quasi-compact scheme, makes \$Y\$ affine.
\item An integral morphism is universally closed, and the image of a universally closed morphism is separated. Morphisms of Spaces, labels spaces-morphisms-lemma-integral-universally-closed (invoked at proof line 3152) and spaces-morphisms-lemma-image-universally-closed-separated (invoked at proof line 3155). They transfer the separatedness of the affine scheme \$Y\$ to the algebraic space \$X\$.
\end{itemize}\medskip
\ifdefined\hypertarget\hypertarget{sourcepointer-1}{}\fi
\noindent\textit{Source reference.} Original lemma 01XF: sources/stacks-chow/tags.txt, source lines 2516--2516. Navigation pointer for \textbackslash{}ref\{coherent-lemma-quasi-compact-h1-zero-covering\} used inside the counted span (line 3149); the argument that consumes it is retained in the span.\par\medskip
\ifdefined\hypertarget\hypertarget{sourcepointer-2}{}\fi
\noindent\textit{Source reference.} Original lemma 09YB: sources/stacks-chow/tags.txt, source lines 12559--12559. Navigation pointer for \textbackslash{}ref\{decent-spaces-lemma-there-is-a-scheme-integral-over\} used inside the counted span (line 3139); the argument that consumes it is retained in the span.\par\medskip
\ifdefined\hypertarget\hypertarget{sourcepointer-3}{}\fi
\noindent\textit{Source reference.} Original lemma 07V0: sources/stacks-spaces-cohomology/pinned-original-source.tex, source lines 3041--3041. Navigation pointer for \textbackslash{}ref\{lemma-vanishing-base-change\} used inside the counted span (line 3142); the argument that consumes it is retained in the span.\par\medskip
\ifdefined\hypertarget\hypertarget{sourcepointer-4}{}\fi
\noindent\textit{Source reference.} Original reference 07UY: sources/stacks-spaces-cohomology/pinned-original-source.tex, source lines 2991--2991. Navigation pointer for \textbackslash{}ref\{situation-vanishing\} used inside the counted span (line 3133); the argument that consumes it is retained in the span.\par\medskip
\ifdefined\hypertarget\hypertarget{sourcepointer-5}{}\fi
\noindent\textit{Source reference.} Original lemma 05Z2: sources/stacks-chow/tags.txt, source lines 7676--7676. Navigation pointer for \textbackslash{}ref\{spaces-morphisms-lemma-image-universally-closed-separated\} used inside the counted span (line 3155); the argument that consumes it is retained in the span.\par\medskip
\ifdefined\hypertarget\hypertarget{sourcepointer-6}{}\fi
\noindent\textit{Source reference.} Original lemma 0415: sources/stacks-chow/tags.txt, source lines 5280--5280. Navigation pointer for \textbackslash{}ref\{spaces-morphisms-lemma-integral-universally-closed\} used inside the counted span (line 3153); the argument that consumes it is retained in the span.\par\medskip
\begin{lemma}
\label{lemma-vanishing-separated}
In Situation \ref{situation-vanishing} the morphism $p : X \to \Spec(A)$ is
separated.
\end{lemma}

\begin{proof}
By Decent Spaces, Lemma
\ref{decent-spaces-lemma-there-is-a-scheme-integral-over}
we can find a scheme $Y$ and a surjective integral morphism
$Y \to X$. Since an integral morphism is affine, we can apply
Lemma \ref{lemma-vanishing-base-change}
to see that $H^1(Y, \mathcal{G}) = 0$ for every
quasi-coherent $\mathcal{O}_Y$-module $\mathcal{G}$.
Since $Y \to X$ is quasi-compact and $X$ is quasi-compact,
we see that $Y$ is quasi-compact.
Since $Y$ is a scheme, we may apply
Cohomology of Schemes, Lemma
\ref{coherent-lemma-quasi-compact-h1-zero-covering}
to see that $Y$ is affine. Hence $Y$ is separated.
Note that an integral morphism is affine and universally closed, see
Morphisms of Spaces, Lemma
\ref{spaces-morphisms-lemma-integral-universally-closed}.
By Morphisms of Spaces, Lemma
\ref{spaces-morphisms-lemma-image-universally-closed-separated}
we see that $X$ is a separated algebraic space.
\end{proof}

\end{document}
